% =========================================================
\section{Introduction}
\label{sec:introduction}
% =========================================================
%
% BEFORE: Knowledge for operation and maintenance (O\&M) in energy-storage stations is inherently dynamic and process-oriented. Handling a single anomaly often requires jointly consulting equipment states, alarms and inspection records, historical work orders, operating procedures, and post-action verification results. Unlike static knowledge question answering, the system must answer two stricter questions in addition to identifying which evidence is relevant to the query: whether that evidence was already available at the query time, and whether the evidence can jointly form a complete process that supports the current decision.
In embodied operation and maintenance (O\&M) for energy-storage stations, decision making depends on dynamic and process-oriented knowledge. Handling a single anomaly often requires jointly consulting equipment states, alarms and inspection records, historical work orders, operating procedures, and post-action verification results. Unlike static knowledge question answering, the system must answer two stricter questions in addition to identifying which evidence is relevant to the query: whether that evidence was already available at the query time, and whether the evidence can jointly form a complete process that supports the current decision.
% 第一二段之间需要增加衔接——第一段是提出现状，第二段介绍RAG，中间需要过度：现状需要RAG来解决
These requirements make retrieval-based grounding useful for O\&M generation because the generator must consult external records that change over time rather than rely only on knowledge stored in model parameters.
Retrieval-augmented generation (RAG) combines language models with external retrieval for knowledge-intensive generation~\cite{lewis2020rag}. On the retrieval side, Best Matching 25 (BM25) remains a strong lexical baseline, while neural retrievers such as DPR and ColBERT improve semantic matching through dense representations or late interaction~\cite{karpukhin2020dpr,khattab2020colbert}; BEIR further shows that retrieval paradigms differ substantially in cross-domain generalization~\cite{thakur2021beir}. On the generation side, FiD demonstrates that a generator can effectively aggregate multiple retrieved passages~\cite{izacard2021fid}, while Self-RAG further studies joint control over retrieval, generation, and self-reflection~\cite{asai2024selfrag}.
%
However, conventional RAG typically treats candidate documents as independent relevance units. This assumption is insufficient for continuously updated O\&M records. First, O\&M evidence has dual temporal attributes: the event time and the time at which the record becomes available to the system need not coincide, and late-arriving records must not be used retrospectively by an earlier query. Second, maintenance decisions usually depend on complementary evidence: a state observation motivates an inspection, the inspection supports a diagnosis, the diagnosis leads to an action, and a post-action check verifies the preceding decision. High scores for individual documents therefore do not guarantee that the Top-$k$ set covers the complete process.

Existing work addresses related aspects of this problem. Time-aware RAG models temporal constraints or information changes across time intervals~\cite{zhang2025mrag,lau2025tarag}; multi-hop retrieval conditions later retrieval steps on previously retrieved information~\cite{xiong2021mdr,trivedi2023ircot}; and RAPTOR, GraphRAG, and HippoRAG organize information across documents using hierarchical or graph structures~\cite{sarthi2024raptor,edge2024graphrag,gutierrez2024hipporag}. These lines of work address temporal filtering, iterative retrieval, and cross-document organization from different angles.
% 需补充类似说明： 然而，上述的方法没有联合解决时间可见性与证据集完整性，因此不直接适用于储能领域
None of the cited methods explicitly combines query-time evidence visibility with direct set-level optimization for O\&M process completeness. This gap motivates a retrieval objective that treats both requirements explicitly.
We propose TEF-RAG (Temporal Evidence Flow Retrieval-Augmented Generation). Its core idea is to construct query-conditioned directed semantic relations among evidence that is visible at the query time, and to reformulate retrieval from independent node ranking into evidence-set ranking. The system first removes unusable records with deterministic asset, temporal, and procedure constraints; it then learns which candidate evidence pairs are worth relation reasoning and constructs a temporal evidence graph; finally, it applies nonlinear pairwise ranking over size-bounded candidate sets so that the model learns which candidate set better satisfies the annotated completeness objective for the current query.

TEF-RAG differs from general graph retrieval in that its graph is not predefined by static entity relations, but is formed dynamically from the currently visible O\&M evidence under the query. It differs from time-aware ranking in that temporal constraints define only the legality boundary, while the final objective remains the recovery of a role-complementary and relation-consistent evidence set. It also differs from stepwise multi-hop retrieval because the target structure may include branching, updates, supersession, and preserved uncertainty rather than being restricted to a single linear path.
% embodied这个描述，如果要有，就全篇要一致，intro第一段有，这里又没有了
% 如果做了embodied指令生成，建议加上，并且介绍为什么需要RAG时也带上相关的描述
In the embodied O\&M setting considered in this work, retrieval is not the endpoint. The selected evidence flow is passed to a shared generator that produces a standardized work order and a dependency-aware action plan, providing an evidence-grounded interface for downstream planners, constraint checkers, and robotic execution. This paper focuses on evidence retrieval and the structured planning context in this pipeline, and does not claim to solve closed-loop robotic control or field safety certification.

Our contributions are threefold:
\begin{enumerate}
    \item We formulate Temporal Evidence Flow retrieval for dynamic O\&M knowledge, unifying temporal visibility and evidence-set completeness within a single retrieval framework.
    \item We propose a method that combines learned evidence-pair proposal, query-conditioned relation graphs, and nonlinear set ranking, shifting the retrieval objective from independent document relevance to set-level process utility.
    % BEFORE: \item We construct a controlled temporal O\&M benchmark and evaluate whether evidence-flow quality transfers to downstream work orders and action plans using both retrieval metrics and structured generation metrics under a shared generator.
    \item We construct a controlled Temporal O\&M benchmark and evaluate retrieval together with downstream structured generation under a shared generator, using both retrieval and evidence-grounded generation metrics.
\end{enumerate}

% =========================================================
\subsection{Related Work}
\label{sec:related_work}
% =========================================================
% Neural Retrieval 似乎没有特意讲到？
\textbf{Retrieval-Augmented Generation and Neural Retrieval.}
Because TEF-RAG operates between candidate retrieval and downstream generation, both neural retrieval and multi-evidence generation provide direct points of comparison.
Classical RAG combines an external retriever with a generator~\cite{lewis2020rag}; DPR and ColBERT represent dense dual-encoder retrieval and late-interaction retrieval, respectively~\cite{karpukhin2020dpr,khattab2020colbert}. FiD demonstrates multi-evidence aggregation from the generation side~\cite{izacard2021fid}, while BEIR systematically compares the generalization of sparse, dense, and reranking paradigms~\cite{thakur2021beir}. DPR and ColBERT rank passages individually, while FiD aggregates multiple retrieved passages at generation time. Our retrieval objective instead assigns utility to an evidence set and its temporal and semantic relations.

% 时间/多跳之间有无关系？ 类似方法/目的相同的方法介绍的时候最好要有一句话连接起来
\textbf{Time-Aware and Multi-Hop Retrieval.}
Time-aware and multi-hop retrieval address different constraints on multi-evidence retrieval: the former models when evidence is relevant, whereas the latter models how multiple pieces of evidence are gathered across steps.
MRAG and TA-RAG explicitly incorporate temporal conditions into retrieval~\cite{zhang2025mrag,lau2025tarag}; MDR and IRCoT find multi-piece evidence through stepwise retrieval or interleaved retrieval and reasoning~\cite{xiong2021mdr,trivedi2023ircot}. TEF-RAG similarly emphasizes complementary evidence, but further distinguishes event time from availability time and allows evidence structures to contain updates, supersession, verification, and uncertainty propagation.

% 同上，另可以加一句引导性的话，不同RAG之间的关系
\textbf{Structured and Graph-Based RAG.}
A complementary line of work changes how retrieved knowledge is organized across documents rather than focusing primarily on temporal filtering or iterative retrieval.
RAPTOR, GraphRAG, and HippoRAG organize cross-document information using hierarchical or graph-based structures~\cite{sarthi2024raptor,edge2024graphrag,gutierrez2024hipporag}. Unlike methods that typically construct graphs from document hierarchies or entity relations, TEF-RAG defines edges jointly from the current query and the currently visible O\&M records; their semantics directly correspond to support, update, qualification, and verification relations among maintenance evidence.

% embodied 要有就都要有
\textbf{Embodied O\&M and Battery Intelligence.}
At the application level, RAG has also begun to support embodied field robotics and professional equipment operation. SayComply studies retrieval of operational-compliance knowledge for field robots, while RAGLRO applies retrieval-augmented language models to robotic operation tasks~\cite{ginting2025saycomply,wang2026raglro}. In the battery domain, LiBrain explores combinations of time-series information, retrieved knowledge, and language models for battery diagnosis and management~\cite{li2026librain}. These studies connect retrieved operational knowledge or time-series information to robotic or battery-management tasks. Our work focuses on the preceding evidence-selection step so that downstream work orders and action plans use evidence that is visible at the query time and covers the required process roles.

% 需要过渡，前面都是RAG系统，这里介绍到L2R
The preceding methods determine what evidence is retrieved or how it is organized; TEF-RAG additionally requires a final mechanism for comparing alternative evidence sets.
\textbf{Learning to Rank.}
The final selection stage of TEF-RAG is a query-level learning-to-rank problem. RankNet learns relative ranking functions from pairwise preferences~\cite{burges2005ranknet}, while work such as ANCE demonstrates the importance of hard negatives for learning effective retrieval boundaries~\cite{xiong2021ance}. We apply this idea to evidence sets rather than individual documents: the training objective directly compares complete sets with high-scoring but incomplete hard negatives under the same query.
